% Pour toute question sur l'utilisation de cette classe, écrivez à jerome.soucy(at)ulaval.ca
% Nécessite la classe Evaluation.sty dans le même dossier ou dans le PATH de votre système
% Nécessite le fichier logo_ul.pdf dans le même dossier un dans le PATH de votre système
\def\Ponderation{35}
\def\SigleNumero{STT-2100}
\def\NomCours{Régression}
\def\DateEvaluation{19 décembre 2023}
\def\DateRemise{}
\def\HeureRemise{}
\def\HeureEvaluation{8h30 à 10h20}
\def\Session{Été 2021}
\def\NomEvaluation{Examen 2}

% Ne rien mettre entre les accolades si l'enseignant (ou l'enseignante) est aussi le coordonnateur (ou la coordonnatrice)
\def\Coordonnatrice{}
\def\Coordonnateur{}

% Ne rien mettre entre les accolades s'il y a plus d'une section
\def\Enseignant{Julien Miron}
\def\Enseignante{}

% Ne rien mettre entre les accolades s'il s'agit d'un cours à section unique
% Mettre le nom de chacun des enseignants et enseignantes entre les accolades
% Une case à cocher sera générée pour chacune des sections
\def\SectionA{}
\def\SectionB{}
\def\SectionC{}
\def\SectionD{}


\def\Directives{
\begin{enumerate}
\item Matériel autorisé:
\begin{itemize}
\item Calculatrice scientifique {\bf autorisée par la FSG};
\item Le feuillet de formules fourni avec l'examen. Celui-ci inclut la table de la loi normale. Veuillez ne rien écrire sur le feuillet de formules.
\end{itemize}
\item Assurez-vous que cette évaluation comporte \numquestions ~exercices répartis sur \numpages{}~pages, pour un total de \numpoints{}~points.\
\item \textbf{Veuillez ne pas dégrafer votre examen.}
\item Ajoutez vos initiales dans le bas de chaque page de l'examen.
\item Veuillez déposer votre carte d'identité \textbf{admise} avec photo sur votre bureau.\
\item Vous devez justifier chaque réponse, soit par une démarche, soit par une (des) phrase(s).
\item  Le verso des feuilles peut être utilisé comme brouillon, \textbf{mais ne sera pas corrigé}, à moins d'une indication claire au recto.
\item Si rien n'est mentionné par rapport aux quatre postulats de RLS et RLM, c'est que ceux-ci sont respectés.
\end{enumerate}
}

\documentclass[Examen]{Evaluation}
\begin{document}
.
\newpage
\begin{questions}  
\question 
Pour chaque énoncé qui suit, indiquez s'il est vrai ou faux.
\textbf{Justifiez} brièvement votre réponse. Il peut s'agir d'une justification mathématique ou d'une phrase. \\
\textit{Aucun point ne sera accordé pour une réponse sans justification ou une mauvaise justification.}

\vspace{0.3cm}
\begin{parts}

\part[4] La méthode de sélection de modèle pas-à-pas avec le critère BIC permet de trouver le modèle avec le plus petit BIC parmi tous les modèles possibles.

     \begin{oneparcheckboxes}
		\choice VRAI
		\correctchoice FAUX
\end{oneparcheckboxes}

\vspace{3 cm}

\part[4] Dans un modèle de régression linéaire multiple, la valeur p du test F ajusté associé à une variable explicative est toujours plus grande que celle du test t partiel de cette même variable.

     \begin{oneparcheckboxes}
		\choice VRAI
		\correctchoice FAUX
\end{oneparcheckboxes}
\vspace{3 cm}

\part[4] Si la valeur-p associée au test-t du coefficient $\beta_j$ d'une certaine variable est de $0,0023$, alors l'hypothèse nulle $H_0:\beta_j=0$ doit être rejetée, lorsque toutes les autres variables sont incluses dans le modèle.

     \begin{oneparcheckboxes}
		\choice VRAI
		\correctchoice FAUX
\end{oneparcheckboxes}
\vspace{3 cm}

\part[4] Si notre modèle de régression inclut l'ordonnée à l'origine, deux variables continues et un terme d'interaction entre ces deux variables, le graphique de l'espérance de $\boldsymbol{Y}$ sera représenté par deux droites.

     \begin{oneparcheckboxes}
		\choice VRAI
		\correctchoice FAUX
\end{oneparcheckboxes}
\end{parts}

\newpage
.
\newpage
\question
\vspace{0.5cm}
\begin{parts}
\part[4] Expliquez pourquoi il pourrait être préférable d'utiliser la distance de Cook plutôt que $DFBETA$.
\vspace{3.5 cm}
\part[4] Marie veut modéliser le temps requis pour que le fromage de sa fondue fonde en fonction de différentes variables explicatives. 
Elle dispose de 6 variables explicatives. Elle applique une sélection de modèle par exclusion basée sur le $R^2$ ajusté. À la fin de la première étape, combien de variables peut-elle avoir dans le modèle? Pour chaque réponse différente, expliquez comment on peut y arriver.

\vspace{3.5 cm}
\part[4] Le $DFFIT$ d'une observation est de -1,7. Que signifie cette valeur?
\vspace{4cm}

\part[4] Avec $n=30$, nous faisons un test global au seuil $\alpha$ pour confronter $H_0:\beta_1=\beta_2=\beta_3=0$ et $H_1:\beta_i\neq0$ pour au moins un $i$, et nous obtenons la statistique $\kappa_{obs}$. Tracez un graphique qui représente la situation où nous rejetons $H_0$. Il s'agit ici de tracer (rapidement) la densité de $\kappa_{obs}$ - \textit{identifiez bien sa loi!} - ainsi que d'ajouter la statistique observée du test et la valeur critique (ou la zone de rejet).

\end{parts}
\vspace{5cm}

\newpage
.
\newpage
\question Avec $i=1,\hdots,n$, nous voulons expliquer une variable réponse $Y_i$ à l'aide de $w_i$, une variable polytomique pouvant prendre les valeurs $\{A,B,C\}$, et la variable quantitative $x_i$. Le modèle de régression utilisé est 

\begin{equation*}
    Y_i=\beta_0 + \beta_1 x_i + \beta_{A} a_{i} + \beta_B b_i + \beta_{IA}a_i x_i+ \beta_{IB}b_i x_i + \varepsilon_i,
\end{equation*}
où 
\begin{align*}
    a_i =\left\{
			\begin{array}{ll}
			1  & \mbox{ si } w_i=A, \\
			0 & \mbox{ sinon, }
   \end{array}
			\right. & & b_i =\left\{
			\begin{array}{ll}
			1  & \mbox{ si } w_i=B, \\
			0 & \mbox{ sinon. }
   \end{array}
			\right. 
\end{align*}

Voici des énoncés basés sur ce modèle:

\begin{enumerate}
    \item Les trois catégories de la variable $w_i$ sont équivalentes pour prédire la variable $Y_i$, peu importe la valeur de la variable $x_i$;
    \item L'effet d'une augmentation de la variable $x_i$ est le même, peu importe la catégorie de la variable $w_i$;
    \item La variable $x_i$ n'a aucun effet sur la variable réponse;
    \item Il n'y a aucun lien linéaire entre la variable réponse et les variables explicatives (ni leur intéraction); 
    \item L'espérance de $Y_i$ est de $0$, lorsque $x_i=0$ et $a_1=1$.
\end{enumerate}

\vspace{0.25cm}
\begin{parts}
    
\part[10] Pour chacun des énoncés ci-haut, cochez l'hypothèse qui correspond.
\begin{table}[h]
\begin{center}
\begin{tabular}{|c||ccccc|}
\hline
& \multicolumn{5}{c|}{Énoncés}\\ \hline
Hypothèses & \multicolumn{1}{l||}{~~~1~~~} & \multicolumn{1}{l||}{~~~2~~~} & \multicolumn{1}{l||}{~~~3~~~} & \multicolumn{1}{l||}{~~~4~~~} & \multicolumn{1}{l||}{~~~5~~~} \\ \hline
$\beta_A=\beta_B$ & \multicolumn{1}{l||}{}  & \multicolumn{1}{l||}{}  & \multicolumn{1}{l||}{}  & \multicolumn{1}{l||}{}  & \multicolumn{1}{l||}{}   \\ \hline
$\beta_A=\beta_B=\beta_{IA}=\beta_{IB}$ & \multicolumn{1}{l||}{}  & \multicolumn{1}{l||}{}  & \multicolumn{1}{l||}{}  & \multicolumn{1}{l||}{}  & \multicolumn{1}{l||}{}     \\ \hline
$\beta_{IA}=\beta_{IB}=0$ & \multicolumn{1}{l||}{}  & \multicolumn{1}{l||}{}  & \multicolumn{1}{l||}{}  & \multicolumn{1}{l||}{}  & \multicolumn{1}{l||}{}     \\ \hline
$\beta_1=\beta_{IA}=\beta_{IB}=0$ & \multicolumn{1}{l||}{}  & \multicolumn{1}{l||}{}  & \multicolumn{1}{l||}{}  & \multicolumn{1}{l||}{}  & \multicolumn{1}{l||}{}     \\ \hline
$\beta_A=\beta_B=\beta_{IA}=\beta_{IB}=0$ & \multicolumn{1}{l||}{}  & \multicolumn{1}{l||}{}  & \multicolumn{1}{l||}{}  & \multicolumn{1}{l||}{}  & \multicolumn{1}{l||}{}     \\ \hline
$\beta_1-\beta_{IA}-\beta_{IB}=0$ & \multicolumn{1}{l||}{}  & \multicolumn{1}{l||}{}  & \multicolumn{1}{l||}{}  & \multicolumn{1}{l||}{}  & \multicolumn{1}{l||}{}     \\ \hline
$\beta_{1}=0$ & \multicolumn{1}{l||}{}  & \multicolumn{1}{l||}{}  & \multicolumn{1}{l||}{}  & \multicolumn{1}{l||}{}  & \multicolumn{1}{l||}{}     \\ \hline
$\beta_1=\beta_A=\beta_B=\beta_{IA}=\beta_{IB}=0$ & \multicolumn{1}{l||}{}  & \multicolumn{1}{l||}{}  & \multicolumn{1}{l||}{}  & \multicolumn{1}{l||}{}  & \multicolumn{1}{l||}{}     \\ \hline
$\beta_0=\beta_1=\beta_A=\beta_B=\beta_{IA}=\beta_{IB}=0$ & \multicolumn{1}{l||}{}  & \multicolumn{1}{l||}{}  & \multicolumn{1}{l||}{}  & \multicolumn{1}{l||}{}  & \multicolumn{1}{l||}{}     \\ \hline
$\beta_0+\beta_A=0$ & \multicolumn{1}{l||}{}  & \multicolumn{1}{l||}{}  & \multicolumn{1}{l||}{}  & \multicolumn{1}{l||}{}  & \multicolumn{1}{l||}{}     \\ \hline
\end{tabular}
\end{center}
\end{table}

\part[6] Faites un graphique de l‘espérance de la variable $Y$ en fonction de $w$ et $x$.
Notez que le graphique n'a pas besoin d'être parfait, ni à l'échelle.
Par contre, on doit y retrouver les éléments importants et le nom des axes.

\end{parts}

\newpage
.
\newpage
\question\label{cook} Soit le modèle de régression

$$\boldsymbol{Y}=\boldsymbol{X\beta}+\boldsymbol{\varepsilon},$$

avec $\boldsymbol{X}=(\boldsymbol{1}, \boldsymbol{x})$ une matrice $n\times 2$, où $\boldsymbol{1}$ est le vecteur de 1 associé à l'ordonnée à l'origine.
Avec un jeu de données de taille $n=31$, on obtient les informations suivantes pour les $3$ premières observations du jeu de données:
\begin{verbatim}
> Sommaire
         x            y       y_pred      Cook
1 28.14979     238.7740     151.0805     0.182
2 26.38498     131.9421     145.2267     0.005
3 28.47468     163.0612     152.1581     0.003
> H[1:3, 1:3]
       [,1]   [,2]   [,3]
[1,] 0.0335 0.0354 0.0331
[2,] 0.0354 0.0402 0.0345
[3,] 0.0331 0.0345 0.0329
\end{verbatim}

Dans ces résultats, $\texttt{y\_pred}$ correspond à $\boldsymbol{\widehat{y}}$, \texttt{Cook} à la distance de Cook et $\texttt{H[1:3, 1:3]}$ aux trois premières lignes et trois premières colonnes de la matrice de projection $\boldsymbol{H}$.

\vspace{0.5cm}
\begin{parts}
    \part[3] Que vaut $\widehat{\sigma}^2$?

    \vspace{4.7cm}
$\widehat{\sigma}^2=$ \boite{3}{1.5}
    \part[3] Que vaut $\widehat{\sigma}_{-1}^2$, la variance estimée sans l'observation 1 ?

\vspace{4.62cm}
$\widehat{\sigma}_{-1}^2=$ \boite{3}{1.5}
\vfill
\pagesuivante{\ref{cook}}
\newpage
.
\newpage
\suite{\ref{cook}}
\part[4] À la lumière de ces résultats et sachant que $p=1$, l'observation 1 a un impact important sur l'estimation de $\beta_1$.

     \begin{oneparcheckboxes}
		\choice VRAI
		\correctchoice FAUX
\end{oneparcheckboxes}

\vspace{4cm}
\end{parts}

\question\label{matricielle}
On examine la réponse d'une variable $Y$ en fonction de deux variables explicatives $v$ et $w$, à partir d'un échantillon de taille $n=30$. Soit $\boldsymbol{X}=(\boldsymbol{1}, \boldsymbol{v}, \boldsymbol{w})$ une matrice $n\times 3$, où $\boldsymbol{1}$ est le vecteur de 1 associé à l'ordonnée à l'origine. Dans les résultats suivants, certaines valeurs de $\boldsymbol{X}^T\boldsymbol{X}$ ont été remplacées par les lettres $a, b, c$ et $d$. Pour les points (c) et (d), $\beta_v$ et $\beta_w$ sont les coefficients associés aux variables $\boldsymbol{v}$ et $\boldsymbol{w}$.
\bigskip

\begin{center}
    $\boldsymbol{X}^T\boldsymbol{X}=\left[
\begin{array}{ccc}
  a &20 & 0\\
 b & 20& 0\\
  c &d & 10\\
\end{array}\right];\qquad
(\boldsymbol{X}^T\boldsymbol{X})^{-1}=\left[
\begin{array}{ccc}
  0,1 &-0,1 & 0\\
 -0,1 & 0,15& 0\\
  0 &0 & 0,1\\
\end{array}\right];\quad
\boldsymbol{X}^T\boldsymbol{Y}=\begin{array}{c}
\end{array}\left[
\begin{array}{r}
  15 \\
  20 \\
 10\\
\end{array}\right];\quad
\boldsymbol{Y}^T\boldsymbol{Y}=59,5
$
\end{center}

\begin{parts}
    \part[4] Que valent $a, b, c$ et $d$? 

    \vspace{4.5cm}

    $a=$\boite{3}{1.5}
    $b=$\boite{3}{1.5}
    $c=$\boite{3}{1.5}
    $d=$\boite{3}{1.5}

\part[6] Que valent $\overline{Y}$, $\overline{v}$ et $\overline{w}$ les moyennes échantillonnales de $\boldsymbol{Y}$, $\boldsymbol{v}$ et $\boldsymbol{w}$?

\vspace{4.5cm}
$\overline{Y}=$\boite{3}{1.5}
$\overline{v}=$\boite{3}{1.5}
$\overline{w}=$\boite{3}{1.5}

\pagesuivante{\ref{matricielle}}
\newpage
.
\newpage
\suite{\ref{matricielle}}
\part[4] En supposant que $\widehat{\sigma}^2=1$, que vaut $\widehat{\text{Var}}(\widehat{\beta}_v)$, l'estimateur de la variance de $\widehat{\beta}_v$?

\vspace{8cm}
$\widehat{\text{Var}}(\widehat{\beta}_1)=$\boite{3}{1.5}

\vspace{0.5cm}
\part[3] En supposant que $\widehat{\sigma}^2=1$, que vaut $\widehat{\text{Cov}}(\widehat{\beta}_v, \widehat{\beta}_w)$, l'estimateur de la covariance entre $\widehat{\beta}_v$ et $\widehat{\beta}_w$?

\vspace{8cm}
$\widehat{\text{Cov}}(\widehat{\beta}_v, \widehat{\beta}_w)=$\boite{3}{1.5}
\end{parts}

\newpage
.
\newpage
\question[10]

Soit le vrai modèle génératif pour $\boldsymbol{Y}$ 
\begin{equation}\label{truemod}
    \boldsymbol{Y} = \boldsymbol{X\beta_1} + \boldsymbol{W\beta_2} + \boldsymbol{Z\beta_3} +\boldsymbol{\varepsilon}.
\end{equation}
Si le modèle ajusté est
\begin{equation}\label{adjmod}
    \boldsymbol{Y} = \boldsymbol{X\beta_1} +\boldsymbol{\varepsilon},
\end{equation}
et que $\widehat{\boldsymbol{\beta}}_1=(\boldsymbol{X}^T\boldsymbol{X})^{-1}\boldsymbol{X}^T\boldsymbol{Y}$ est utilisé pour estimer $\boldsymbol{\beta}_1$ à partir du modèle \eqref{adjmod}, calculez le biais de $\widehat{\boldsymbol{\beta}}_1$ pour $\boldsymbol{\beta}_1$.
\newpage
.
\newpage
\question\label{analyse}
On s'intéresse à la qualité de différents vins (mesurée par une note). La sortie \texttt{R} représente la régression linéaire effectuée.

 Utilisez cette sortie \texttt{R}  pour répondre aux questions suivantes. {\bf Justifiez toutes vos réponses.}
\begin{verbatim}
Coefficients:
                       Estimate   Std. Error   t value   Pr(>|t|)    
(Intercept)             21.9700      21.1900     1.036     0.3002    
acidite.fixe             0.0250       0.0260     0.963     0.3357    
acidite.volatire        -1.0840       0.1211    -8.948      2e-15
acide.citrice           -0.1826       0.1472    -1.240     0.2150    
sucre.residuel           0.0163       0.0150     1.089     0.2765    
chlorures               -1.8740       0.4193    -4.470   8.37e-06 
anhydrire.sulf.libre     0.0044       0.0022     2.009     0.0447 
dioxyde.soufre.total    -0.0033       0.0007    -4.480   8.00e-06 
densite                -17.8800      21.6300    -0.827     0.4086    
pH                      -0.4137       0.1916    -2.159     0.0310 
sulfates                 0.9163       0.1143     8.014   2.13e-15 
alcool                   0.2762       0.2648    10.429    < 2e-16 

Residual standard error:   ?   on  1587  degrees of freedom
Multiple R-squared:    ? ,	Adjusted R-squared:   0.3561 
F-statistic:  81.35  on  11  and  1587  DF,  p-value:  < 2.2e-16    
\end{verbatim}

\begin{parts}
    \part[2]  Combien de vins font partie de l'échantillon?
    \vspace{1in}
    \part[4] La valeur-p associée au test sur la variable \texttt{alcool} est la plus petite parmi toutes les variables. 
    Aurélien affirme qu'il est donc certain que la variable \texttt{alcool} est celle qui expliquerait le mieux la variable réponse, dans une régression linéaire simple. Commentez.

    \vfill

    \pagesuivante{\ref{analyse}}
    \newpage
    .
\newpage
    \suite{\ref{analyse}}
    \part[4] Calculez le pourcentage de la variabilité de la qualité du vin qui est expliquée par le modèle.
    \vspace{8cm}


    \part[5] Lorsque Jérôme magasine son vin, il aime choisir les vins en fonction de leur qualité. Il aimerait choisir le modèle qui prédit le mieux possible la qualité du vin en fonction des variables explicatives disponibles. Après quelques manipulations, il obtient les résultats suivants:
	\begin{center}\begin{tabular}{|c|c|c|c|c|c|}\hline
		Mod\`ele & $R^2$ & $R_{a}^2$ & AIC & BIC & $PRESS$\\ \hline\hline
		Modèle 1 & 0.3606 & 0.3561 & 3164.277 & 3234.179 & 679.2157 \\\hline
            Modèle 2 & 0.3603 & 0.3562 & 3162.965 & 3227.490 & 678.3656 \\\hline
            Modèle 3 & 0.3602 & 0.3565 & 3161.224 & 3220.372 & 677.3349 \\\hline
            Modèle 4 & 0.3599 & 0.3567 & 3158.839 & 3207.610 & 676.3202 \\\hline
            Modèle 5 & 0.3595 & 0.3566 & 3159.977 & 3213.371 & 675.8162 \\\hline
            Modèle 6 & 0.3594 & 0.3527 & 3173.829 & 3216.847 & 681.8370 \\\hline
	\end{tabular}\end{center}
    Sur la base de ce tableau, quel modèle est préférable pour la mise en situation décrite? Justifiez votre réponse, un choix de modèle non justifié ne recevra aucun points.
    
    \vspace{2in}

\end{parts}


%%%%%%%




 
	



\end{questions}
\newpage
.
\newpage



\end{document}




\vspace{0.3in}
\textbf{Question 4 (10 pts)} (Design d'expérience)\\

On s'intéresse à un médicament sensé réduire le nombre de crises chez les individus atteints d'epilepsie. Vous devez aider à planifier une expérience pour étudier l'effet de la dose du médicament sur le nombre de crises. Puisque vous vous intéressez à la relation entre deux variables quantitatives, vous planifiez analyser les données avec un modèle de régression linéaire simple. \\

Soit $x_1, \ldots, x_n$ les valeurs de la dose pour les 20 individus de votre étude.\\ Pour chacun, vous pouvez attribuer n'importe quelle dose entre 2 mg et 21 mg. \\

\begin{itemize}
%\item [(a)] Pour le modèle de régression linéaire simple $Y_i = \beta_0 + \beta_1 x_i + \varepsilon_i$, prouvez que la variance de $\beta_0$ est minimale si les $x_i$ sont tels que $\bar{x}$ = 0.
\item [(a)] (5 pts) Quelle est la meilleure fa\c{c}on de choisir les valeurs $x_1,\ldots,x_n$ pour minimiser la variance de votre estimateur $\hat{\beta_1}$? Justifiez votre réponse.
\vspace{0.1 in}
\item [(b)] (5 pts) Pourquoi ne serait-ce pas vraiment une bonne idée de choisir les valeurs de $x_i$ de cette façon, même si cela minimise la variance de votre estimateur?
\end{itemize}

\end{document} 